\chapter{研究工作进度安排}
\section{排期依据和主要依赖}
研究排期围绕资料、金标和真实输入组织，而不是按模型名称依次运行。收益标签和共享四折已经提供底座，当前最先影响后续训练的是名单/主体/历史版本证据与人工预告金标。日装配器可在CPU先建立接口和边界测试，但真实Qwen缓存需要抽取模型评价和冻结后才生成。这些依赖决定了哪些任务能够并行，哪些只能等待前置条件。

\reportfigure{来源证据、人工标注、CPU接口与正式实验的依赖}{fig:dependencies}{\FigDependencies}{研究实施设计；箭头表示产物依赖而非已完成比例。来源核验和人工标注可与CPU接口并行，正式训练在前置条件满足后启动。}

图\ref{fig:dependencies}中的并行部分具有现实意义。等待数据库回复时，可进行真实窗口、事件版本流、缺失掩码及索引检查；等待标注裁决时，可完成普通门控和强树的信息接口。另一方面，随机文本向量或未认证边不能用于正式比较，不能为了赶排期先跑一轮再把这些替代输入写成实证。

\section{目标阶段和可检查产物}
表\ref{tab:schedule}使用目标时间，不代替学院校历。学校两年制专业学位硕士开题应在第三学期第1周（含）前完成，用户入学年级和正式考核安排尚未提供，因此本稿的编制日期不能证明该手续按时完成。研究工作安排由导师结合个人培养计划复核。

\reporttable{阶段工作与产物安排}{tab:schedule}{
\begin{tabularx}{\linewidth}{p{2.45cm}YY}\toprule
目标阶段&主要工作&提交或核查产物\\\midrule
2026年11月底&补充开题正文、图表与文献；定位来源证据&导师可审稿、文献阅读记录、原报告任务表\\
12月上旬&名单/主体/证券核验，200/42独立试标；CPU真实接口&原版问题分类、裁决及规范版本、窗口/时点测试\\
12月中旬&完成剩余1500/300标注；共同信息对照和恢复调度&人工放行数据、SFT整批预检、日装配与配置\\
前置条件满足后&GPU长度/显存/吞吐试跑，抽取器独立评价和缓存&冻结模型与缓存键、资源日志、真实输入核查\\
目标12月下旬至1月&四折双方五种子主比较、固定消融及P5&验证选择台账、全部种子结果、配对统计\\
目标2027年2月&整合论文、复核引用数字、导师修改与答辩&完整论文及非敏感复现材料；正式截止依学院\\\bottomrule
\end{tabularx}}

人工标注约1800次判定加裁决，工时应从200条试标记录估计。300双标是主样本子集，不额外增加300篇公告。正式全矩阵还包含对照验证配置、各种子重训、消融和P5，不能只按“一个模型一轮”计算时间。GPU小时、I/O等待和中断恢复分别计入，再预留运行余量。

\section{排期偏离时怎样调整}
若历史版本无法取得，首先列清无法认证的组、所占年份和覆盖变化，明确可接受与隔离范围；不把最新修正版自动当作历史原版。若金标分歧较多，先修订训练试标规范并补裁决，而不是提前把不一致答案送入SFT。若真实吞吐低于预期，先检查窗口重复、邻居重复编码、缓存I/O和恢复粒度，再调整计算批次及租赁时段。

这些调整改变的是工程执行和资源安排，不应改变共同测试样本、主指标或测试选模规则。方案保留完整三路编码、关系消息和有界融合，实验结果可以否定部分假设；不能在看到失败结果后删掉弱年份、种子或对照，再把剩余结果称作预先确定的实验。

阶段成果还应与培养手续区分。导师内容审核、政治立场和价值导向审核、开题口头考核及系统上传需要本人和培养单位实际办理。文稿、程序和实验产物为这些环节提供材料，不自动产生审核意见、签名或考核结果。
